%!TEX TS-program = xelatex
%!TEX encoding = UTF-8
% 管理世界（Management World）格式模板
% 参考来源：http://www.mgmt.gov.cn/
% 适用于管理学、经济学、社会学等领域实证论文
% CTeX 中文格式，使用XeLaTeX编译

\PassOptionsToPackage{quiet}{xeCJK}
\documentclass[lang=cn, aaffiliation={中国}]{ctexart}

% ── 文档类型选项 ────────────────────────────────────────────────────────────
% 管理世界采用单栏排版，12pt（五号）字体，1.5 倍行距

% ── 必要宏包 ────────────────────────────────────────────────────────────────
\usepackage[
    a4paper,
    top=2.54cm, bottom=2.54cm,
    left=3.17cm, right=3.17cm,
    headheight=14pt,
    heightrounded,
]{geometry}

\usepackage{fontspec}
\usepackage{xeCJK}
\setmainfont{Times New Roman}
\setsansfont{Arial}
\xeCJKsetup{CJKspace=true, xCJKecglue={}}

\usepackage{mathtools}
\usepackage{booktabs}
\usepackage{threeparttable}
\usepackage{longtable}
\usepackage{graphicx}
\graphicspath{{figures/}}
\usepackage{hyperref}
\hypersetup{
    colorlinks=true,
    linkcolor=blue,
    citecolor=blue,
    urlcolor=blue,
    pdftitle={管理世界格式模板},
    pdfauthor={作者姓名},
}

\usepackage{natbib}
\bibliographystyle{gbt7714-numerical}

\usepackage{setspace}
\usepackage{indentfirst}
\usepackage{enumitem}
\usepackage{pdflscape}          % 横向表格
\usepackage{afterpage}          % 浮动体位置控制

% ── 页面设置 ────────────────────────────────────────────────────────────────
% 管理世界要求：五号（10.5pt）字体，1.5 倍行距
\setstretch{1.5}
\setlength{\parindent}{2em}
\linespread{1.5}

% ── 章节标题 ────────────────────────────────────────────────────────────────
\titleformat*{\section}{\fontsize{12pt}{1.5em}\bfseries\center}
\titleformat*{\subsection}{\fontsize{12pt}{1.5em}\bfseries}
\titleformat*{\subsubsection}{\fontsize{12pt}{1.5em}\rmfamily}

% ── 页眉页脚 ────────────────────────────────────────────────────────────────
\usepackage{fancyhdr}
\pagestyle{fancy}
\fancyhf{}
\fancyhead[C]{\footnotesize 管~理~世~界}
\fancyfoot[C]{\thepage}

% ── 数学命令 ────────────────────────────────────────────────────────────────
\DeclareMathOperator{\E}{E}
\DeclareMathOperator{\Var}{Var}
\DeclareMathOperator{\Cov}{Cov}

% ── 标题与作者 ─────────────────────────────────────────────────────────────
\title{论文标题（中文）\thanks{基金项目（编号）；基金项目（编号）。}}
\author{作者一\thanks{简介：\dots，电话：\dots，E-mail：\dots}
       \quad 作者二\thanks{简介：\dots，E-mail：\dots}}
\date{\today}

% ── 摘要 ────────────────────────────────────────────────────────────────────
\newenvironment{cnabstract}{%
    \par\textbf{摘\quad 要}
    \indent
}{\par\vspace{0.5\baselineskip}\par\textbf{关键词：}...}

% ═══════════════════════════════════════════════════════════════════════════
% 正文
% ═══════════════════════════════════════════════════════════════════════════

\begin{document}

\thispagestyle{empty}
\maketitle
\thispagestyle{fancy}

% ── 摘要 ────────────────────────────────────────────────────────────────────
\begin{cnabstract}
    本文基于\dots数据，研究了\dots的问题。实证结果表明：（1）\dots；
    （2）\dots；（3）\dots。进一步分析发现，\dots是重要的影响机制。
    本文的研究发现对\dots具有重要的政策启示。
\end{cnabstract}

\textbf{关键词：}关键词1；关键词2；关键词3；关键词4；关键词5

% ── 一、引言 ───────────────────────────────────────────────────────────────
\section{一、引言}
\label{sec:intro}

\indent 在当前\dots背景下，\dots问题日益受到学术界和实践界的关注。

从理论层面看，\dots（\citealp{xxx}）。\citet{yyy}的研究指出，\dots。

从实践层面看，\dots（\citealp{zzz}）。

然而，已有研究仍存在以下不足：第一，\dots；第二，\dots；第三，\dots。

针对上述研究缺口，本文利用\dots数据，系统考察了\dots的影响及其机制。

本文可能的边际贡献在于：（1）\dots；（2）\dots；（3）\dots。

% ── 二、文献综述 ──────────────────────────────────────────────────────────
\section{二、文献综述}
\label{sec:lit}

\subsection{（一）\dots研究}
\label{sec:lit:topic1}

\citet{xxx}最早系统研究了\dots问题。

后续研究主要从以下两个视角展开：

\textbf{视角一：}基于\dots理论，\dots（\citealp{yyy}）。\citet{zzz}进一步发现\dots。

\textbf{视角二：}基于\dots框架，\citealp{www}认为\dots。

\subsection{（二）\dots研究}
\label{sec:lit:topic2}

关于\dots的研究主要集中于\dots。

\citet{vvv}利用\dots数据，发现\dots。

综合而言，现有文献存在以下局限：（1）\dots；（2）\dots；（3）\dots。

\subsection{（三）研究假说推导}
\label{sec:hypothesis}

基于上述文献梳理，本文提出如下研究假说：

\textbf{假说 H1：}在其他条件不变的情况下，\dots对\dots有显著的正向（负向）影响。

\textbf{假说 H2：}\dots对\dots的影响存在异质性，具体表现为\dots。

\textbf{假说 H3：}\dots在\dots对\dots的影响中发挥中介作用。

% ── 三、研究设计 ──────────────────────────────────────────────────────────
\section{三、研究设计}
\label{sec:design}

\subsection{（一）样本选择与数据来源}
\label{sec:data}

本文以\dots为研究对象，时间跨度为\dots年。

数据来源：（1）\dots数据库；（2）\dots数据库；（3）\dots官方统计。

样本筛选过程：（1）剔除\dots样本；（2）剔除\dots样本；（3）对\dots进行处理。

最终得到包含\dots个观测值的面板数据集。

\subsection{（二）变量定义}
\label{sec:var}

表~\ref{tab:var_def}~详细定义了本文使用的主要变量。

\begin{table}[htbp]
    \centering
    \caption{主要变量定义}
    \label{tab:var_def}
    \begin{threeparttable}
        \begin{tabular}{p{2.2cm}p{2.0cm}p{6.3cm}}
            \toprule
            \textbf{变量类别} & \textbf{变量名称} & \textbf{变量定义与说明} \\
            \midrule
            \textbf{因变量} & $Y_{it}$ & \dots \\
            \midrule
            \textbf{自变量} & $X_{it}$ & \dots \\
              & $Z_{it}$ & \dots \\
            \midrule
            \textbf{中介变量} & $M_{it}$ & \dots \\
            \midrule
            \textbf{调节变量} & $W_{it}$ & \dots \\
            \midrule
            \textbf{控制变量} & Size & \dots \\
              & Lev & \dots \\
              & Age & \dots \\
              & $ROA$ & \dots \\
              & $Dual$ & \dots \\
              & $Top1$ & \dots \\
            \midrule
            \textbf{固定效应} & $\mu_i$ & 企业固定效应 \\
              & $\lambda_t$ & 年份固定效应 \\
            \bottomrule
        \end{tabular}
        \begin{tablenotes}\footnotesize
            \item（1）所有连续变量均在 1\% 和 99\% 分位数进行 Winsorize 处理；
            （2）数据来源：CSMAR、Wind 及国家统计局。
        \end{tablenotes}
    \end{threeparttable}
\end{table}

\subsection{（三）模型设定}
\label{sec:model}

为检验假说 H1，本文构建如下双向固定效应模型：

\begin{equation}
    \label{eq:baseline}
    Y_{it} = \alpha + \beta_1 X_{it} + \beta_2 Z_{it} + \gamma Controls_{it} + \mu_i + \lambda_t + \varepsilon_{it},
\end{equation}

其中，$i$ 表示企业，$t$ 表示年份，$\mu_i$ 和 $\lambda_t$ 分别为企业和年份固定效应，
用于控制\dots。

为检验假说 H3（中介效应），本文采用三步法：

\begin{align}
    \label{eq:mediation}
    M_{it} &= \alpha_1 + c \cdot X_{it} + \gamma_1 Controls_{it} + \mu_i + \lambda_t + \varepsilon_{it}, \\
    Y_{it} &= \alpha_2 + a \cdot X_{it} + b \cdot M_{it} + \gamma_2 Controls_{it} + \mu_i + \lambda_t + \varepsilon_{it}.
\end{align}

% ── 四、实证分析 ──────────────────────────────────────────────────────────
\section{四、实证分析}
\label{sec:results}

\subsection{（一）描述性统计}
\label{sec:desc}

表~\ref{tab:desc}~报告了主要变量的描述性统计结果。

\begin{longtable}{lccccccc}
    \caption{描述性统计}
    \label{tab:desc}\\
    \toprule
    \textbf{变量} & \textbf{N} & \textbf{Mean} & \textbf{SD} & \textbf{Min} & \textbf{P25} & \textbf{P75} & \textbf{Max} \\
    \midrule
    \endfirsthead
    \caption{(续表~\ref{tab:desc})}\\
    \toprule
    \textbf{变量} & \textbf{N} & \textbf{Mean} & \textbf{SD} & \textbf{Min} & \textbf{P25} & \textbf{P75} & \textbf{Max} \\
    \midrule
    \endhead
    \bottomrule
    \endfoot
    $Y_{it}$ & 12,480 & 0.052 & 0.098 & -0.215 & -0.004 & 0.085 & 0.302 \\
    $X_{it}$ & 12,480 & 0.318 & 0.466 & 0.000 & 0.000 & 1.000 & 1.000 \\
    $M_{it}$ & 12,480 & 0.184 & 0.388 & 0.000 & 0.000 & 0.000 & 1.000 \\
    Size & 12,480 & 22.16 & 1.308 & 19.72 & 21.19 & 22.81 & 25.62 \\
    Lev & 12,480 & 0.412 & 0.192 & 0.051 & 0.246 & 0.548 & 0.856 \\
    Age & 12,480 & 2.87 & 0.302 & 2.197 & 2.708 & 3.091 & 3.466 \\
\end{longtable}

从表~\ref{tab:desc}~可以发现：（1）\dots；（2）\dots；（3）\dots。

各主要变量之间的相关系数（未汇报，备索）均处于合理范围，不存在严重的多重共线性问题。

\subsection{（二）基准回归结果}
\label{sec:baseline}

表~\ref{tab:regression}~报告了基准回归结果。

\begin{table}[htbp]
    \centering
    \caption{基准回归结果}
    \label{tab:regression}
    \begin{threeparttable}
        \begin{tabular}{lcccc}
            \toprule
            \textbf{} & \textbf{(1)} & \textbf{(2)} & \textbf{(3)} & \textbf{(4)} \\
            \midrule
            $X_{it}$ & 0.038\*** & 0.028\*** & 0.022\** & 0.019\** \\
              & (4.12) & (3.28) & (2.47) & (2.34) \\
            $Z_{it}$ & & 0.015\** & 0.012\** & 0.010\** \\
              & & (2.58) & (2.31) & (2.18) \\
            \midrule
            Size & & & -0.003 & -0.002 \\
              & & & (-0.68) & (-0.54) \\
            Lev & & & -0.021\*** & -0.019\*** \\
              & & & (-3.85) & (-3.52) \\
            ROA & & & 0.128\*** & 0.116\*** \\
              & & & (5.42) & (4.89) \\
            \midrule
            \textbf{固定效应} & \textbf{无} & \textbf{无} & \textbf{无} & \textbf{企业+年份} \\
            \textbf{Adj. $R^2$} & 0.054 & 0.071 & 0.098 & 0.186 \\
            \textbf{N} & 12,480 & 12,480 & 12,480 & 12,480 \\
            \bottomrule
        \end{tabular}
        \begin{tablenotes}\footnotesize
            \item 注：（1）括号内为 $t$ 值，标准误在企业层面进行聚类调整；
            （2）\*** $p<0.01$，\** $p<0.05$，\* $p<0.10$；
            （3）所有回归均控制了常数项。
        \end{tablenotes}
    \end{threeparttable}
\end{table}

由表~\ref{tab:regression}~可见：（1）\dots；（2）\dots；（3）\dots。

\subsection{（三）平行趋势检验}
\label{sec:parallel}

为保证 DID 估计的无偏性，本文采用事件研究法（Event Study）检验平行趋势假设。

图~\ref{fig:parallel}~的结果显示，政策实施前各期系数均不显著（$p>0.10$），
表明处理组与对照组在政策实施前满足平行趋势假设。

\begin{figure}[htbp]
    \centering
    \caption{平行趋势检验：事件研究图}
    \label{fig:parallel}
    % \includegraphics[width=0.9\textwidth]{event_study.pdf}
    \begin{minipage}{0.9\textwidth}\footnotesize
        注：横轴为相对时间（0 表示政策实施年份），纵轴为 DID 系数估计值，
        虚线表示 95\% 置信区间。政策实施前的系数均在统计意义上与零无显著差异，
        支持平行趋势假设。
    \end{minipage}
\end{figure}

% ── 五、进一步分析 ────────────────────────────────────────────────────────
\section{五、进一步分析}
\label{sec:additional}

\subsection{（一）异质性分析}
\label{sec:hetero}

表~\ref{tab:hetero}~报告了分样本回归结果。

\begin{table}[htbp]
    \centering
    \caption{异质性分析}
    \label{tab:hetero}
    \begin{threeparttable}
        \begin{tabular}{lcccc}
            \toprule
            \textbf{} & \textbf{(1)} & \textbf{(2)} & \textbf{(3)} & \textbf{(4)} \\
            & \textbf{国有企业} & \textbf{非国有企业} & \textbf{东部地区} & \textbf{中西部地区} \\
            \midrule
            $X_{it}$ & 0.012 & 0.028\*** & 0.024\*** & 0.008 \\
              & (1.23) & (3.52) & (3.14) & (0.86) \\
            \midrule
            \textbf{Adj. $R^2$} & 0.168 & 0.194 & 0.178 & 0.202 \\
            \textbf{N} & 3,840 & 8,640 & 7,280 & 5,200 \\
            \bottomrule
        \end{tabular}
        \begin{tablenotes}\footnotesize
            \item 注：所有控制变量与固定效应均与基准回归一致；\*** $p<0.01$，\** $p<0.05$。
        \end{tablenotes}
    \end{threeparttable}
\end{table}

\subsection{（二）稳健性检验}
\label{sec:robust}

为确保研究结论的可靠性，本文进行了以下稳健性检验：

\textbf{（1）替换核心变量的度量方式。}
采用\dots作为\dots的替代变量，主要结论保持不变。

\textbf{（2）改变样本范围。}
剔除\dots样本后重新回归，结果仍然稳健。

\textbf{（3）调整控制变量集。}
参考\citet{xxx}的经典研究，加入\dots等变量，结果保持一致。

\textbf{（4）改变标准误聚类层级。}
同时在行业和年份层面进行双维聚类，结果依然稳健。

% ── 六、结论与启示 ────────────────────────────────────────────────────────
\section{六、结论与启示}
\label{sec:conclusion}

本文以\dots为准自然实验，利用\dots数据，系统研究了\dots的影响及其机制。

\textbf{主要结论：}
第一，\dots；第二，\dots；第三，\dots。

\textbf{政策启示：}
第一，\dots；第二，\dots；第三，\dots。

\textbf{参考文献：}
\bibliography{references}

\end{document}
